The ``ExpandConnects'' pass expands the connections between aggregate
types into those between ground types. It also flattens the elements
in aggregate types to their corresponding ground types in declaration
statements. To expand the objects of aggregate types in a Firrtl
module, the process involves two phases. The first phase is to record
the types along with orientations and flipping information of the
objects, which helps determine the actual connection directions. For
example, if there is an input port of aggregate types with a field
marked as ``flip'', the ``flip'' mark will recursively apply to all the
atomic elements in that field. Since it is an input port, the
orientation is record as ``source''. The second phase is to expand the statements
sequentially. For declaration statements, ground-typed elements are
generated corresponding to the original aggregate types. For
connection statements, the connection directions are determined by the
orientations and flipping information from the first phase. For
subfields, there are two kinds of connections allowed for HiFirrtl
programs. One is to set connections between elements of equivalent
aggregate types. In such bundled connections, individual subfield
connections can be of different directions. The other kind is to
connect subfields separately, in which case the direction should
comply with the orientation of the subfield elements.


At the beginning of a Firrtl module, ports are declared first. When
ports are defined with flip sub-fields, the expansion process will
change those fields of input ports into output ports, and vice
versa. As a result, their orientations are changed between ``source''
and ``sink''.
 \begin{gather*}
   \label{expand-input}
   \infer[]
         {\seqq{{\tt expand\_connect} (\mbox{\prettyll{input v}}) = ss}}
         {\begin{array}{c}
             \seqq{{\tt expand\_ref \;v} = [v_0, v_1,\dots v_n]}\\
             \seqq{ s s= {\tt map}\; (fun\; v_x \Rightarrow if\;{\tt is\_flipped}(v_x)\;then\; \mbox{\prettyll{output}} {\tt v_x}\;else\; \mbox{\prettyll{input}} {\tt v_x})\, [v_0, v_1,\dots v_n]}
           \end{array}
         }
 \end{gather*}

Elements such as wires and registers have duplex
orientation. Regardless of the flip sub-fields in the original
aggregate types, the resulting orientations will remain duplex. To
expand registers in aggregate types, we also need to expand the reset
values of the registers. Expanding wires is relatively simple and is
omitted here.
\begin{gather*}
  \label{expand-reg}
  \infer[]
        {\seqq{{\tt expand\_connect} (\mbox{\prettyll{reg v : t clk with : (reset => rst, rv)}}) = ss}}
        {\begin{array}{c}
            \seqq{{\tt expand\_ref \;}v = [v_0,\dots v_n]}\\
            \seqq{{\tt expand\_expr \;}rv = [rv_0,\dots rv_n]}\\
            \seqq{ s s= {\tt map}\; (fun\; v_x\; rv_x \Rightarrow \mbox{\prettyll{reg v_x: t clk with:(reset=> rst,rv_x)}})}
            %% \rightline{
              \seqq{[v_0,\dots v_n], [rv_0,\dots rv_n]}
          \end{array}
        }
\end{gather*}

Nodes are also ``source'' elements. They are only allowed to be declared
with passive aggregate types, which means no flip sub-fields. During
the expansion of node declarations, the passive type is checked;
otherwise, the result will be an empty sequence.
\begin{gather*}
  \label{expand-node}
  \infer[]
        {\seqq{{\tt expand\_connect} (\mbox{\prettyll{node v e}}) = ss}}
        {\begin{array}{c}
            \seqq{{\tt expand\_ref}\; v = [v_0, v_1,\dots v_n]}\\
            \seqq{{\tt expand\_expr}\; e = [e_0, e_1,\dots e_n]}\\
            \seqq{ s s= if\;{\tt is\_passive}(v_x)\; then\; {\tt map}\; (fun\; v_x\;e_x \Rightarrow \; \mbox{\prettyll{node}} {\tt v_x}\;{\tt e_x})\,[v_0, v_1,\dots v_n]\,[e_0, e_1,\dots e_n] else\; [::]}
          \end{array}
        } 
\end{gather*}

To expand \prettyll{is invalid} statements, we
select only the ``sink'' and ``duplex'' elements to produce the sequence of
invalidations.
In HiFirrtl, it is allowed to have ``source'' sub-elements in a bundle be
invalidated. The result of expansion will omit them.
\begin{gather*}
  \label{expand-invalid}
  \infer[]
        {\seqq{{\tt expand\_connect} (\mbox{\prettyll{v is invalid}}) = ss}}
        {\begin{array}{c}
            \seqq{{\tt expand\_ref \;}v = [v_0, v_1,\dots v_n]}\\
            \seqq{vs = {\tt filter}\; (fun\; v_x \Rightarrow orient (v_x) \neq source) [v_0, v_1,\dots v_n]}\\
            \seqq{ s s = {\tt map}\; (fun\; v_x \Rightarrow {\tt v_x}\;\mbox{\prettyll{is invalid}})\, vs}
          \end{array}
        } 
\end{gather*}

Connections between bundled elements contain a set of connections for
individual sub-fields. According to the indicated ``flip'' for a sub-field, the
connection direction may differ from the original.
\begin{gather*}
   \label{expand-invalid}
   \infer[]
         {\seqq{{\tt expand\_connect} (\mbox{\prettyll{l <= r}}) = ss}}
         {\begin{array}{c}
             \seqq{{\tt expand\_ref }\; l = [l_0,\dots l_n]}\\
             \seqq{{\tt expand\_expr}\; r = [r_0,\dots r_n]}\\
             \seqq{ s s = {\tt map}\; (fun\; l_x\;r_x \Rightarrow if\;{\tt is\_flipped}(l_x)\; then\;\mbox{\prettyll{r_x <= l_x}}\; else\mbox{\prettyll{l_x <= r_x}})\,[l_0,\dots l_n]\,[r_0,\dots r_n]}
           \end{array}
         }
\end{gather*}



